\section{Method}
\label{sec:method}

\Cref{fig:architecture} gives the full picture.
Frames arrive one at a time from a live video stream. At each step $t$, the Semantic Keyframe Memory (\cref{sec:skm}) scores the new frame $I_t$ and decides whether to store it in a bank of at most $N$ entries.
Independently, a user may submit a question $q$ at any moment.
The Temporal Query Router (\cref{sec:tqr}) classifies $q$ into a temporal scope and passes only the matching memory entries forward.
The Stream-Fused Generator (\cref{sec:sfg}) then runs a two-stage perception-synthesis pipeline over those entries and returns an answer in $\sim$250\,ms on an A100 GPU.

\subsection{Semantic Keyframe Memory}
\label{sec:skm}

A live video stream produces frames at 15 to 30 fps. Consecutive frames are redundant.
Storing all frames is not feasible. Uniform subsampling loses rare but informative events.
The SKM retains frames based on visual-semantic importance.

\paragraph{Visual encoding.}
Each incoming frame $I_t$ passes through a CLIP ViT-B/32~\cite{radford2021clip} encoder, producing a 512-dimensional embedding $v_t$.
Not every frame goes through encoding. By default the system processes every 5th frame, giving 3 to 6 candidates per second and keeping encoding cost low.
The system saves a small JPEG thumbnail alongside each embedding for later visualization.

\paragraph{Importance scoring.}
The system scores every candidate against what is already in memory. Two signals determine the score.

\textbf{Visual novelty} captures how different the candidate looks. We compute the cosine similarity between $v_t$ and each stored embedding, and one minus the maximum gives the novelty. If every stored frame looks similar to the candidate, novelty is low. If no stored frame is close, novelty is high.

\textbf{Temporal coverage} captures how well the candidate fills a gap in the timeline. It is the smallest absolute time difference between $t$ and any stored timestamp, divided by a horizon $T_{\max}$.

The importance score for frame $I_t$ is:
\begin{equation}
\begin{split}
  s(I_t) =\; & \alpha \Bigl(1 - \max_{j \in \mathcal{M}} \cos(v_t, v_j)\Bigr) \\
  & + (1{-}\alpha)\,\min\!\Bigl(\frac{\min_{j \in \mathcal{M}} |t - t_j|}{T_{\max}},\, 1\Bigr)
\end{split}
  \label{eq:importance}
\end{equation}
where $\mathcal{M}$ is the current memory set, $\alpha = 0.7$ controls the novelty-coverage trade-off, and $T_{\max} = 300$\,s (5 minutes). A frame that is both visually distinct and temporally distant from stored entries receives the highest score.

\paragraph{Memory update.}
Below capacity, every candidate enters memory directly.
Once memory is full, the system re-scores all stored entries and finds the lowest-scoring one. If the incoming candidate scores higher, it replaces that entry; otherwise the candidate is dropped. The full update runs in $O(N)$ time and adds under 1\,ms of latency.
Because importance is recomputed against the \emph{current} memory state at each insertion, a frame's score evolves over time: a frame that was once unique may become redundant as similar frames arrive later.
This dynamic re-scoring is the key difference from one-pass selection methods such as SeViLA~\cite{yu2023sevila}, which commit to a fixed set once.

\paragraph{Temporal bookkeeping.}
Each memory entry stores its original timestamp alongside a compact JPEG thumbnail (used for the demo filmstrip) and its CLIP embedding (used for novelty computation). The TQR uses these timestamps to filter entries by time window. We sort memory by timestamp after each insertion so that scope-filtered retrieval (\cref{sec:tqr}) requires a single binary search.


\subsection{Temporal Query Router}
\label{sec:tqr}

When a question $q$ arrives, the TQR determines which part of memory is relevant.
The pipeline has three stages: a keyword matcher scores the query against per-scope keyword lists, a scope selector picks the best temporal scope, and a memory filter returns only entries in the matching time window.

\paragraph{Scope classification.}
The TQR assigns each question one of three temporal scopes: \textsc{instant}, \textsc{recent}, or \textsc{historical}.
A keyword matcher runs over the lowercased query string.
Three manually curated keyword lists drive the decision. \textsc{Instant} cues include ``right now,'' ``currently,'' and ``at this moment.'' \textsc{Recent} cues include ``just,'' ``a moment ago,'' ``few seconds,'' and ``last minute.'' \textsc{Historical} cues include ``earlier,'' ``before,'' ``previously,'' and ``how many times.''

The scope with the most keyword hits wins. When no keyword matches or two scopes tie with low confidence (top score / total $<$ 0.6), the system defaults to \textsc{historical}, granting full memory access. This fallback is conservative: an over-broad context is better than a truncated one.

\paragraph{Adapting to event density.}
The three scopes are defined \emph{linguistically}, not temporally: they reflect how the user phrases the question rather than imposing fixed time ranges.
A question like \emph{``What just happened?''} is \textsc{recent} regardless of whether the stream shows one event per minute (surveillance) or ten events per second (sports).
The scope label determines \emph{which slice of memory} to consult; the slice's temporal extent is controlled separately by the recency window~$W$.

In practice, different video streams exhibit very different event densities.
When $W = 30$\,s, the \textsc{recent} scope captures roughly 3--5 events in a high-density cooking stream or 0--1 events in a slow surveillance feed.
For dense streams, the SFG's uniform sampling (\cref{sec:sfg}) selects a representative subset.
For sparse streams where the window may be empty, the conservative default to \textsc{historical} ensures the system searches the full memory rather than returning an empty context.
We chose a single $W$ for simplicity; an adaptive scheme that estimates event density from the SKM insertion rate could refine this boundary automatically, but the fixed value works well across all evaluated domains (\cref{sec:experiments}).

\paragraph{Scope-aware filtering.}
Given predicted scope $c$ and current time $t_{\text{now}}$, the filtered subset $\mathcal{F} \subseteq \mathcal{M}$ is:
\begin{equation}
  \mathcal{F} \!=\!
  \begin{cases}
    \{\argmax_{j \in \mathcal{M}} t_j\}                        & c = \text{inst.} \\[3pt]
    \{j \!\in\! \mathcal{M} : t_j \!\ge\! t_{\text{now}} - W\} & c = \text{rec.}  \\[3pt]
    \mathcal{M}                                                 & c = \text{hist.}
  \end{cases}
  \label{eq:scope_filter}
\end{equation}
where $W = 30$\,s is the recency window. \textsc{inst.}, \textsc{rec.}, and \textsc{hist.}\ stand for the three scopes defined above. Only entries in $\mathcal{F}$ reach the Stream-Fused Generator.


\subsection{Stream-Fused Generator}
\label{sec:sfg}

The SFG turns filtered keyframes into a natural-language answer through two stages. \Cref{alg:sfg} gives the full procedure.

\paragraph{Stage 1: BLIP visual perception.}
The system samples up to $K\!=\!6$ frames from the filtered set $\mathcal{F}$ (or takes all of them if $|\mathcal{F}| \le K$).
When $|\mathcal{F}| > K$, we select uniformly spaced indices to preserve temporal coverage.
For non-instant queries, the most recent frame is always appended to the sampled set, anchoring the answer in the present even when the question targets the past.
Two BLIP~\cite{li2022blip} models process each sampled frame $I_k$ independently:

\begin{itemize}[itemsep=2pt]
  \item \textbf{BLIP captioning} (blip-image-captioning-base): generates a free-form description of the frame contents. Each caption is at most 50 tokens.
  \item \textbf{BLIP VQA} (blip-vqa-base): takes the user's question $q$ and the frame $I_k$ as input. It produces a short direct answer for that specific frame. Each answer is at most 30 tokens.
\end{itemize}

The system collects all captions and VQA answers into a pool of textual observations.
Deduplication removes near-identical entries using word-level Jaccard similarity (threshold $\ge 0.6$, ignoring stop words); the remaining unique observations are ranked by frequency so that observations consistent across multiple frames receive higher weight.
For each frame, the system also runs two targeted VQA probes (``Where is this?'' and ``What is happening?'') and folds the answers into the caption to produce enriched scene descriptions.

\paragraph{Stage 2: Flan-T5 language synthesis.}
The system assembles a scope-specific prompt from the user's question, deduplicated captions, and deduplicated VQA outputs.
For \textsc{instant} queries, the prompt asks for a direct description of the single frame.
For \textsc{recent} and \textsc{historical} queries, the prompt asks Flan-T5 to summarize the observations into a coherent narrative, suppressing generic phrases.
Yes/no questions use a constrained format that forces Flan-T5 to begin with ``Yes'' or ``No.''
Flan-T5-base~\cite{chung2022flanT5} generates up to 100 tokens with beam search (2 beams, length penalty 1.0, no-repeat 3-gram constraint).
If the generation echoes the prompt or is too short ($\le 5$ words), the system falls back to a direct answer assembled from the highest-frequency observations.

\begin{algorithm}[t]
\SetAlgoLined
\KwIn{Question $q$, memory $\mathcal{M}$, time $t_{\text{now}}$}
\KwOut{Answer string $a$}
$c \leftarrow \texttt{TQR}(q)$ \tcp*{scope classification}
$\mathcal{F} \leftarrow \texttt{ScopeFilter}(\mathcal{M}, c, t_{\text{now}})$\;
$S \leftarrow \texttt{Sample}(\mathcal{F}, K\!=\!6)$\;
\If{$c \neq \text{instant}$}{
  $S \leftarrow S \cup \{\argmax_{j \in \mathcal{M}} t_j\}$ \tcp*{append live frame}
}
$\mathcal{C}, \mathcal{V} \leftarrow \emptyset, \emptyset$\;
\ForEach{frame $I_k \in S$}{
  $\mathcal{C} \leftarrow \mathcal{C} \cup \{\texttt{BLIP-Cap}(I_k)\}$\;
  $\mathcal{V} \leftarrow \mathcal{V} \cup \{\texttt{BLIP-VQA}(I_k, q)\}$\;
}
$p \leftarrow \texttt{BuildPrompt}(q, \texttt{dedup}(\mathcal{C}), \texttt{dedup}(\mathcal{V}))$\;
$a \leftarrow \texttt{Flan-T5}(p)$\;
\Return{$a$}
\caption{Stream-Fused Generator}
\label{alg:sfg}
\end{algorithm}

\paragraph{Latency budget.}
The total per-query latency breaks down as follows. TQR scope classification takes under 1\,ms. Frame sampling and BLIP inference on up to 6 frames takes 80--120\,ms with batched FP16 inference. Flan-T5 decoding takes 56\,ms (2 beams).
On an NVIDIA A100 GPU, end-to-end latency averages 203\,ms per query in isolated profiling and 254\,ms across the full evaluation, comfortably supporting real-time interaction (see \cref{sec:latency} for the full breakdown).


\subsection{Training and Hyperparameters}
\label{sec:training}

Every model in the pipeline is a frozen pre-trained checkpoint; no end-to-end training is required.
The vision encoder is CLIP ViT-B/32~\cite{radford2021clip} (224$\times$224 input, 512-d embeddings).
Two BLIP~\cite{li2022blip} checkpoints handle perception: blip-image-captioning-base and blip-vqa-base, both loaded in evaluation mode.
The language synthesis model is Flan-T5-base~\cite{chung2022flanT5}, instruction-tuned on 1,836 tasks.
We selected $\alpha$ and $T_{\max}$ by sweeping $\alpha \in \{0.5, 0.6, 0.7, 0.8\}$ and $T_{\max} \in \{60, 120, 300, 600\}$ on a held-out set, choosing the combination $\alpha = 0.7$, $T_{\max} = 300$\,s by LiveQA-Bench accuracy.
TQR keyword lists were compiled by hand from annotated question-scope pairs in the LiveQA-Bench training split and validated on a held-out split.
